\chapter{Complexity — Space Amplify Poly}
%%% generated by the blueprint pipeline from TCSlib.Complexity.Randomized.SpaceAmplifyPoly
\section{Overview}

Declarations of the Lean module \texttt{TCSlib.Complexity.Randomized.SpaceAmplifyPoly}.

\section{Declarations}

\begin{definition}[Binary length with an added unit]
\lean{Randomized.PolyAmp.ell}
\label{Randomized.PolyAmp.ell}
\leanok
For a natural number $n$, let $\ell(n)$ be one more than the length of its binary digit
list.
\end{definition}

\begin{definition}[Exponential repetition parameter]
\lean{Randomized.PolyAmp.bigK}
\label{Randomized.PolyAmp.bigK}
\leanok
\uses{Randomized.PolyAmp.ell}
For natural numbers $F$ and $n$, define $K=2^{F\ell(n)}$, where $\ell(n)$ is one more
than the length of the binary digit list of $n$.
\end{definition}

\begin{definition}[Initial register configuration]
\lean{Randomized.PolyAmp.sreg}
\label{Randomized.PolyAmp.sreg}
\leanok
\uses{Randomized.PolyAmp.Reg}
For natural numbers $n,l,s,e,K,R,t,a$, define a register configuration with these values
in registers $0$ through $5$, respectively; set registers $6$ through $10$ and register
$14$ to zero, and set registers $11$ and $12$ to $t$ and $a$, respectively. Set register
$13$ to zero.
\end{definition}

\begin{definition}[Zero register configuration]
\lean{Randomized.PolyAmp.reg0}
\label{Randomized.PolyAmp.reg0}
\leanok
\uses{Randomized.PolyAmp.Reg}
Define the register configuration in which all fifteen registers, numbered $0$ through
$14$, are zero.
\end{definition}

\begin{lemma}[Register-vector invariance of probabilistic semantics]
\lean{Randomized.PolyAmp.PS.of_v}
\label{Randomized.PolyAmp.PS.of_v}
\leanok
\uses{Complexity.LogProg.PSem, Complexity.LogProg.RegLe, Randomized.PolyAmp.Lab, Randomized.PolyAmp.PS, Randomized.PolyAmp.Reg, Randomized.PolyAmp.Reg.v, Randomized.PolyAmp.ampA, Randomized.PolyAmp.ampOr}
\difficulty{1}
Let $N$ be a finite nondeterministic Turing machine with Boolean input symbols, let $o$
be a Boolean parameter, let $x$ be a Boolean input word, let $B$ be a natural-number
register bound, let $l$ be a control label, let $p$ be rational, let $\sigma'$ be a
register state, and let $f$ be a function from the fifteen register indices to the
natural numbers. If $f$ is the register vector of $\sigma'$ and the
polynomial-amplification program, started at $l$ with register vector $\sigma'$ and no
result, has every register at most $B$, terminates for every sequence of choice bits,
and returns true and false with probabilities $p$ and $1-p$, respectively, for every
number of choice bits for which every branch terminates, then the abstract program has
the same termination and probability properties when started at $l$ with register vector
$f$ and no result.
\end{lemma}

\begin{theorem}[Initial amplifier semantics]
\lean{Randomized.PolyAmp.ps_start}
\label{Randomized.PolyAmp.ps_start}
\leanok
\uses{Complexity.Codec.halvCount_loop, Complexity.LogProg.AReach.mono, Complexity.LogProg.NARM.detPart, Complexity.LogProg.PSem.inc, Complexity.LogProg.PSem.of_reach, Complexity.LogProg.inputLength_loop, Complexity.LogProg.registerScale_loop, Complexity.LogProg.shl_run, Complexity.indicator_indexLang_pairEncode, Complexity.inputLenLang, Randomized.PolyAmp.Lab, Randomized.PolyAmp.PS, Randomized.PolyAmp.PS.of_v, Randomized.PolyAmp.Reg.Le, Randomized.PolyAmp.Reg.le_v, Randomized.PolyAmp.Reg.v, Randomized.PolyAmp.ampA, Randomized.PolyAmp.ampAs, Randomized.PolyAmp.ampOr, Randomized.PolyAmp.bigK, Randomized.PolyAmp.codeBound, Randomized.PolyAmp.cpDst, Randomized.PolyAmp.cpSrc, Randomized.PolyAmp.ell, Randomized.PolyAmp.mk, Randomized.PolyAmp.ps_cp, Randomized.PolyAmp.ps_runs, Randomized.PolyAmp.reg0, Randomized.PolyAmp.runVal, Randomized.PolyAmp.sL0, Randomized.PolyAmp.sL1, Randomized.PolyAmp.sreg, Turing.FinNDTM.UsesSpace}
\difficulty{5}
Let $N$ be a finite-state nondeterministic Turing machine with Boolean input, let $x$ be
a Boolean input word, and let $sp:\mathbb{N}\to\mathbb{N}$ bound the space used by $N$
on every input and every choice word. For $n\in\mathbb{N}$, write $\ell(n)$ for one plus
the length of the binary digit list of $n$, and set $K_*=2^{F\ell(|x|)}$. Suppose
$sp(|x|)\le cS\,\ell(|x|)$, the configuration-code bound for $N,|x|,cS\,\ell(|x|)$ is at
most $K_*$, and $|x|+\ell(|x|)+(cS+F)\ell(|x|)+2K_*+2\le B$. Then the amplifier started
at its initial label with all registers zero terminates for every sequence of choice
bits, keeps each register at most $B$ while running, and has acceptance probability
equal to the run value with parameters $K_*$, $K_*+1$, and $0$.
\end{theorem}

\begin{theorem}[Lower Bound for the Run Value]
\lean{Randomized.PolyAmp.runVal_two_ge}
\label{Randomized.PolyAmp.runVal_two_ge}
\leanok
\uses{Randomized.PolyAmp.accBit, Randomized.PolyAmp.decRule, Randomized.PolyAmp.runVal, Randomized.blockCount, Randomized.blockCount_add_not, Randomized.randProb, Randomized.randProb_bool_false, Randomized.randProb_congr, Randomized.randProb_mono, Randomized.randProb_not, Randomized.randProb_tail_le}
\difficulty{5}
Let $N$ be a finite-state nondeterministic Turing machine over Boolean symbols and let
$x$ be a Boolean input. For any nonnegative integer $K$, if the probability that $N$
accepts on a uniformly random choice word of length $K$ is at least $2/3$, then the run
value with the rejecting decision rule, $K$ blocks, initial count $0$, and parameter
$i=K+1$ is at least $1-2(8/9)^{\left\lfloor (K+1)/2 \right\rfloor}$.
\end{theorem}

\begin{theorem}[Upper Bound for the Run Value]
\lean{Randomized.PolyAmp.runVal_two_le}
\label{Randomized.PolyAmp.runVal_two_le}
\leanok
\uses{Randomized.PolyAmp.accBit, Randomized.PolyAmp.decRule, Randomized.PolyAmp.runVal, Randomized.blockCount, Randomized.randProb, Randomized.randProb_mono, Randomized.randProb_tail_le}
\difficulty{4}
Let $N$ be a finite-state nondeterministic Turing machine over Boolean symbols and let
$x$ be a Boolean input. For any nonnegative integer $K$, if the probability that $N$
accepts on a uniformly random choice word of length $K$ is at most $1/3$, then the run
value with the rejecting decision rule, $K$ blocks, initial count $0$, and parameter
$i=K+1$ is at most $2(8/9)^{\left\lfloor (K+1)/2 \right\rfloor}$.
\end{theorem}

\begin{theorem}[Acceptance probability after repeated trials]
\lean{Randomized.PolyAmp.runVal_one}
\label{Randomized.PolyAmp.runVal_one}
\leanok
\uses{Randomized.PolyAmp.accBit, Randomized.PolyAmp.decRule, Randomized.PolyAmp.runVal, Randomized.randProb, Randomized.randProb_blockCount, Randomized.randProb_congr, Randomized.randProb_not}
\difficulty{4}
For every finite-state nondeterministic Turing machine $N$, Boolean input $x$, and
$K\in\mathbb N$, the run value with the one-sided decision rule, $K$ trials per block,
$K+1$ blocks, and initial count zero equals $1-(1-p)^{K+1}$, where $p$ is the fraction
of uniformly random length-$K$ choice words on which $N$ accepts from its initial
configuration on $x$.
\end{theorem}

\begin{theorem}[Exponential Bound from a Size Condition]
\lean{Randomized.PolyAmp.two_mul_pow_le}
\label{Randomized.PolyAmp.two_mul_pow_le}
\leanok
\difficulty{3}
For nonnegative integers $K,n,d$, if $32(n+1)^d\le K$, then $2(8/9)^{\left\lfloor (K+1)/2 \right\rfloor}\le
(1/2)^{(n+1)^d}$.
\end{theorem}

\begin{theorem}[Bound on the Binary Digit Length]
\lean{Randomized.PolyAmp.two_pow_bits_le}
\label{Randomized.PolyAmp.two_pow_bits_le}
\leanok
\difficulty{5}
For every nonnegative integer $n$, the number of binary strings of length equal to the
binary digit length of $n$ is at most $2n+1$; equivalently, $2^{\ell(n)}\le 2n+1$, where
$\ell(n)$ is the length of the binary expansion of $n$ (with length zero when $n=0$).
\end{theorem}

\begin{theorem}[Binary bit-length bound]
\lean{Randomized.PolyAmp.lt_two_pow_bits}
\label{Randomized.PolyAmp.lt_two_pow_bits}
\leanok
\difficulty{1}
For every natural number $n$, one has $n < 2^{\lvert\operatorname{bits}(n)\rvert}$,
where $\operatorname{bits}(n)$ is the binary expansion of $n$.
\end{theorem}

\begin{theorem}[Logarithmic length comparison]
\lean{Randomized.PolyAmp.logSpace_le_ell}
\label{Randomized.PolyAmp.logSpace_le_ell}
\leanok
\uses{Complexity.logSpace, Randomized.PolyAmp.ell, Randomized.PolyAmp.lt_two_pow_bits}
\difficulty{4}
For every natural number $n$, the quantity $\lfloor\log_2 n\rfloor+1$ is at most one
more than the length of the binary representation of $n$.
\end{theorem}

\begin{theorem}[Lower bound for the amplification parameter]
\lean{Randomized.PolyAmp.bigK_ge}
\label{Randomized.PolyAmp.bigK_ge}
\leanok
\uses{Randomized.PolyAmp.bigK, Randomized.PolyAmp.ell, Randomized.PolyAmp.lt_two_pow_bits}
\difficulty{5}
For all natural numbers $F,n,d$, if $d+5\le F$, then $32(n+1)^d\le 2^{F\ell(n)}$, where
$\ell(n)$ is one more than the length of the binary representation of $n$.
\end{theorem}

\begin{theorem}[Upper bound for the amplification parameter]
\lean{Randomized.PolyAmp.bigK_le}
\label{Randomized.PolyAmp.bigK_le}
\leanok
\uses{Randomized.PolyAmp.bigK, Randomized.PolyAmp.ell, Randomized.PolyAmp.two_pow_bits_le}
\difficulty{4}
For all natural numbers $F,n$, $2^{F\ell(n)}\le 4^F(n+1)^F$, where $\ell(n)$ is one more
than the length of the binary representation of $n$.
\end{theorem}

\begin{theorem}[Binary representation length bound]
\lean{Randomized.PolyAmp.ell_le}
\label{Randomized.PolyAmp.ell_le}
\leanok
\uses{Randomized.PolyAmp.ell, Randomized.PolyAmp.lt_two_pow_bits}
\difficulty{3}
For every natural number $n$, the length of its binary representation plus one is at
most $n+1$.
\end{theorem}

\begin{theorem}[Polynomial register bound]
\lean{Randomized.PolyAmp.regBound_le}
\label{Randomized.PolyAmp.regBound_le}
\leanok
\uses{Randomized.PolyAmp.bigK, Randomized.PolyAmp.bigK_le, Randomized.PolyAmp.ell, Randomized.PolyAmp.ell_le}
\difficulty{4}
For all natural numbers $cS$, $F$, and $n$, if $\ell(n)$ is one plus the length of the
binary representation of $n$ and $K_*=2^{F\ell(n)}$, then
$n+\ell(n)+(cS+F)\ell(n)+2K_*+2\le(cS+F+6+2\cdot4^F)(n+1)^{F+1}$.
\end{theorem}

\begin{theorem}[Polynomial amplification setting]
\lean{Randomized.polyAmp_setting}
\label{Randomized.polyAmp_setting}
\leanok
\uses{Complexity.Codec.codeK, Complexity.LogProg.NHalts, Complexity.LogProg.NHalts.mono, Complexity.LogProg.NPreS, Complexity.LogProg.PSem.decides, Complexity.LogProg.ansProb, Complexity.logSpace, Randomized.PolyAmp.Reg.v, Randomized.PolyAmp.accBit, Randomized.PolyAmp.accP_init, Randomized.PolyAmp.ampA, Randomized.PolyAmp.ampA_npreS, Randomized.PolyAmp.ampOr, Randomized.PolyAmp.bigK, Randomized.PolyAmp.codeBound, Randomized.PolyAmp.ell, Randomized.PolyAmp.haltsWithin_codeBound, Randomized.PolyAmp.logSpace_le_ell, Randomized.PolyAmp.ps_start, Randomized.PolyAmp.reg0, Randomized.PolyAmp.regBound_le, Randomized.PolyAmp.runVal, Randomized.randProb, Randomized.randProb_congr, Turing.FinNDTM, Turing.FinNDTM.UsesSpace, Turing.FinNDTM.outputProb, Turing.FinNDTM.outputProb_eq_of_haltsWithin, Turing.NDTM.HaltsWithin, Turing.NDTM.HaltsWithin.mono}
\difficulty{5}
Let $N$ be a finite nondeterministic Turing machine with Boolean input symbols, and let
$c,F$ be natural numbers. Fix a Boolean parameter $os$, and suppose that on every input
and every choice word, $N$ visits at most $c$ times the logarithmic space bound for the
input length, and that the configuration-code width of $N$, multiplied by $c+2$, is at
most $F$. For a Boolean input word $x$, suppose every computation of $N$ on $x$ halts
within $T_0$ steps, and set $K=2^{F\ell(|x|)}$. Then every computation on $x$ halts
within $K$ steps, and for every output word the probability of that output after $K$
random bits equals its probability after $T_0$ random bits. The probability that a
uniformly random $K$-bit choice word makes $N$ accept equals the probability that its
output is the one-bit word $[\mathrm{true}]$. Moreover, there is a natural number $N_c$
such that the amplification program halts after $N_c$ random bits and its register
values, at each live step before halting, are bounded by
$(c+F+6+2\cdot4^F)(|x|+1)^{F+1}$; its probabilities of returning true and false are
respectively the value obtained by applying the stated decision rule to the block count
in the amplification experiment and one minus that value.
\end{theorem}

\begin{theorem}[Error reduction for bounded-error log space]
\lean{Randomized.BPL_error_reduction}
\label{Randomized.BPL_error_reduction}
\leanok
\uses{Complexity.Codec.codeK, Complexity.LogProg.narm_logspace_ptm, Complexity.logSpace, Randomized.BPL, Randomized.PolyAmp.accBit, Randomized.PolyAmp.ampA, Randomized.PolyAmp.ampAs, Randomized.PolyAmp.ampAs_mem, Randomized.PolyAmp.bigK, Randomized.PolyAmp.bigK_ge, Randomized.PolyAmp.runVal_two_ge, Randomized.PolyAmp.runVal_two_le, Randomized.PolyAmp.two_mul_pow_le, Randomized.polyAmp_setting, Randomized.randProb, Randomized.randProb_le_one, Randomized.randProb_or_disjoint, Turing.FinNDTM, Turing.FinNDTM.UsesSpace, Turing.FinNDTM.outputProb, Turing.MultiTapeTM.indicator, Turing.NDTM.HaltsWithin}
\difficulty{5}
Let $L$ be a language over the Boolean alphabet that belongs to bounded-error
probabilistic logarithmic space, and let $d$ be a natural number. There exist a natural
number $c$ and a finite-state nondeterministic Turing machine $M$ over the Boolean
alphabet such that, for every input $x$ and every sequence of nondeterministic choices,
$M$ uses at most $c(\lfloor\log_2 |x|\rfloor+1)$ work-tape cells. Moreover, for every
input $x$ there is a natural number $T$ such that every choice sequence of length $T$
makes $M$ halt on $x$, and the probability that its output is the membership bit of $x$
in $L$ is at least $1-(1/2)^{(|x|+1)^d}$.
\end{theorem}

\begin{theorem}[Error Reduction for Randomized Logspace]
\lean{Randomized.RL_error_reduction}
\label{Randomized.RL_error_reduction}
\leanok
\uses{Complexity.Codec.codeK, Complexity.LogProg.narm_logspace_ptm, Complexity.logSpace, Randomized.PolyAmp.Accepting, Randomized.PolyAmp.accBit, Randomized.PolyAmp.ampA, Randomized.PolyAmp.ampAs, Randomized.PolyAmp.ampAs_mem, Randomized.PolyAmp.bigK, Randomized.PolyAmp.bigK_ge, Randomized.PolyAmp.runVal_one, Randomized.RL, Randomized.exists_ofFn_eq, Randomized.forall_of_randProb_eq_one, Randomized.polyAmp_setting, Randomized.randProb, Randomized.randProb_congr, Randomized.randProb_false, Randomized.randProb_le_one, Turing.FinNDTM, Turing.FinNDTM.UsesSpace, Turing.FinNDTM.outputProb, Turing.FinNDTM.outputWith, Turing.FinNDTM.runWith_eq_runWith_take, Turing.NDTM.HaltsWithin, Turing.NDTM.runWith_append, Turing.NDTM.runWith_of_halt}
\difficulty{6}
Let $L$ be a language of Boolean strings in randomized logspace, and let $d$ be a
nonnegative integer. There exist a constant $c$ and a finite-state randomized Turing
machine $M$ using at most $c(\lfloor\log_2 |x|\rfloor+1)$ space on inputs $x$ such that,
for every $x$, there is a time bound $T$ by which every computation branch halts. If
$x\in L$, then $M$ outputs the one-symbol string $[\mathrm{true}]$ with probability at
least $1-(1/2)^{(|x|+1)^d}$; if $x\notin L$, then every choice word of length $T$ makes
$M$ output $[\mathrm{false}]$.
\end{theorem}
